\section{Conclusion}

This research proposal wishes to explore the relationship between media coverage and government advertising spending in Mexico, focusing on how a sharp reduction in spending during Andrés Manuel López Obrador's (AMLO) presidency influenced media behavior. Historically, Mexican governments have used state resources, including advertising, to influence narratives and suppress dissent. The main hypothesizes is that this significant reduction led to increased critical coverage of the government, particularly on the issue of violence, which remained salient across administrations. To analyze this, a difference-in-differences (DiD) design with continuous treatment will be employed, using changes in spending as the treatment variable and incorporating state-level variations in violence coverage. The dataset will include detailed government advertising records and a near-complete archive of online newspaper articles from 2012 to 2024. Articles will be classified using machine learning to assess their tone (positive or negative) and topic (political or violence-related), providing insights into how financial pressures affect media reporting in a democratic context.